\documentclass[10pt,a4paper]{article}
\usepackage[utf8]{inputenc}
\usepackage{pgfplots}
\pgfplotsset{compat=1.15}
\usepackage{mathrsfs}
\usetikzlibrary{arrows}
\pagestyle{empty}
\usepackage{amsmath}
\usepackage{amsfonts}
\usepackage{amssymb}
\usepackage{graphicx}
\usepackage[left=2cm,right=2cm,top=2cm,bottom=2cm]{geometry}
\usepackage{listings}
\usepackage{color}
\definecolor{munsell}{rgb}{0.0, 0.5, 0.69}
\definecolor{dkgreen}{rgb}{0,0.6,0}
\definecolor{gray}{rgb}{0.5,0.5,0.5}
\definecolor{mauve}{rgb}{0.58,0,0.82}
\definecolor{minas}{RGB}{0.244, 0.172, 0.36}
\definecolor{PUJ}{RGB}{44, 86, 151}
\definecolor{PUJ2}{RGB}{43, 93, 156}
\definecolor{PUJ3}{RGB}{20, 52, 107}
\definecolor{cyandk}{rgb}{0.0, 0.72, 0.92}

\author{Iván Andrés Trujillo }

%% code information
\lstset{frame=tb,
  language=Python,
  aboveskip=3mm,
  belowskip=3mm,
  showstringspaces=false,
  columns=flexible,
  basicstyle={\small\ttfamily},
  numbers=none,
  classoffset=1,
  morekeywords={True,False}, keywordstyle=\color{munsell}, 
  classoffset=0, 
  keywordstyle=\color{blue},  
  commentstyle=\color{dkgreen},
  stringstyle=\color{PUJ3},
  numberstyle=\tiny\color{gray},
  breaklines=true,
  breakatwhitespace=true,
  tabsize=4,
}

\begin{document}

\newcommand{\HRule}{\rule{\linewidth}{0.5mm}}

%----------------------------------------------------------------------------------------
%	LOGO SECTION
%----------------------------------------------------------------------------------------

\includegraphics[width = 4cm]{pujshield.eps}\\[0.5cm] 

\begin{center} 
%%%% Title%%%%%
\textbf{\LARGE Introduction to heuristics and meta-heuristics  in Python}\\[0.2cm]
\emph{\LARGE Notes of class}\\[0.3cm] 
%%%%%%%%%%%%%%%%%%%%%%
\emph{Iván Andrés Trujillo Abella} \\
\textsc{\Large 
}\\[0.2cm] 
\textsc{\large Facultad de ingeniería}\\[0.5cm] 
\HRule \\[0.4cm]
\end{center}
\vspace{1cm}




\section{Some concepts about graph theory}
a graph $G = (V,E) $ composed of two set the set of vertex $V$ and the set of edges $E$.
the set $V$ could be person, items or another entities while $E$ is the set that linked the them.

for instance:
$V = \{ (A,B,C) \}$ and $E =  \{ (A,B), (A,C), (B,C) \}$.
generally the graphical representation is by points and lines, where vertex and points and lines that connected the vertices represented edges.

in some cases could be exist loops inside graphs for instance:
$V = \{ (A,B,C,D \}$ and $E = \{ (A,A), (A,B),(C,D) \}$ note that so far we had been dispensing of the loop $(A,A)$.


\section{Algorithms and search}

\begin{itemize}
\item algorithm hill climbing
\item search A*
\item minmax 
\item alfa and beta algorithm
\end{itemize}






rubik cube.


peeling: peladura
fading: desvanecimiento.

mainstream:
circunvent:
braid:

amusement:
noughts;
scrambled:
spotlight:
shelves:
halted:
advertisements:
rear: trasero,
craze:
plummeted:
shortage:
sued:
appealed:
overturned:
amended:
disclosed:
upheld:
concealed:
inward:
screw:
advertised:
flipped:
scrambled:
takes up:
lament;
awful:
silinesss



\section{Algorithms in blind}

expaction function 
and stop function:




Depth first search (DFS):
it is an algorithm, but not guaranteed the optimum route.


stack structure:



hash set:




\subsection{DFS algorithm}
depth first search 
how works here the data structure ?


important here longitute of path and branching factor:


\subsubsection{LIFO structure }
LIFO (last in first out ), the stack is a example, the last element added is the first removed. Are useful to implement recursion.

\begin{lstlisting}
stack = ['A','B','C']
stack.append('D') # Add to the tail 'D'
stack.pop() # remove the last element namely 'D'
\end{lstlisting}



\subsection{BFS algorithm}


\subsubsection{FIFO structure}
FIFO (first in first out), the queue is a example, the first element added is the first element removed.
imagine good or services, line in a restaurant.


\begin{lstlisting}
queue  = ['A','B','C','D']
queue.pop(0) # Remove the first element  'A'
queue.append('Z') # Add the 'A' to the tail.
\end{lstlisting}





\section{Applying BFS and DFS to a particular tree}



\definecolor{xdxdff}{rgb}{0.49019607843137253,0.49019607843137253,1}
\begin{tikzpicture}[line cap=round,line join=round,>=triangle 45,x=1cm,y=1cm]
\clip(4.4460283220320145,3.4839563756562133) rectangle (26.028671106012656,17.35447680358096);
\draw [line width=2pt] (13.9,14.74) circle (1cm);
\draw [line width=2pt] (12,12.16) circle (1cm);
\draw [line width=2pt] (16.08,12.36) circle (1cm);
\draw [line width=2pt] (12.6,9.3) circle (1cm);
\draw [line width=2pt] (9.78,9.5) circle (1cm);
\draw [line width=2pt] (16.07875233780117,9.398522089690589) circle (1cm);
\draw [line width=2pt] (18.66,9.66) circle (1cm);
\draw [line width=2pt] (13.303297233655439,13.937537659053866)-- (12.371390676354103,13.08847669088526);
\draw [line width=2pt] (14.378852130680574,13.862104427085615)-- (15.315088801882949,13.004135745782781);
\draw [line width=2pt] (11.385571860132973,11.371027211534171)-- (10.08561234285114,10.452156024975443);
\draw [line width=2pt] (11.781660987974588,11.184127018599364)-- (12.241095078760932,10.233374125155812);
\draw [line width=2pt] (16.1610200714651,11.363287529916581)-- (16.07014909872466,10.39848508114447);
\draw [line width=2pt] (16.63968150395849,11.531292202204684)-- (17.999046964532916,10.41042726823247);
\draw (13.796623559594307,15.10863897461962) node[anchor=north west] {$\mathbf{A}$};
\draw (11.833281055030623,12.566181055040746) node[anchor=north west] {$\mathbf{B}$};
\draw (16.014211856115878,12.74980301589922) node[anchor=north west] {$\mathbf{C}$};
\draw (9.60156799228917,9.938974538142576) node[anchor=north west] {$\mathbf{D}$};
\draw (12.41239647004581,9.656479213744923) node[anchor=north west] {$\mathbf{E}$};
\draw (15.943588025016465,9.882475473263046) node[anchor=north west] {$\mathbf{F}$};
\draw (18.55666977569475,9.938974538142576) node[anchor=north west] {$\mathbf{G}$};
\begin{scriptsize}
\draw [fill=xdxdff] (13.303297233655439,13.937537659053866) circle (2.5pt);
\draw [fill=xdxdff] (12.371390676354103,13.08847669088526) circle (2.5pt);
\draw [fill=xdxdff] (14.378852130680574,13.862104427085615) circle (2.5pt);
\draw [fill=xdxdff] (15.315088801882949,13.004135745782781) circle (2.5pt);
\draw [fill=xdxdff] (11.385571860132973,11.371027211534171) circle (2.5pt);
\draw [fill=xdxdff] (10.08561234285114,10.452156024975443) circle (2.5pt);
\draw [fill=xdxdff] (11.781660987974588,11.184127018599364) circle (2.5pt);
\draw [fill=xdxdff] (12.241095078760932,10.233374125155812) circle (2.5pt);
\draw [fill=xdxdff] (16.1610200714651,11.363287529916581) circle (2.5pt);
\draw [fill=xdxdff] (16.07014909872466,10.39848508114447) circle (2.5pt);
\draw [fill=xdxdff] (16.63968150395849,11.531292202204684) circle (2.5pt);
\draw [fill=xdxdff] (17.999046964532916,10.41042726823247) circle (2.5pt);
\end{scriptsize}
\end{tikzpicture}







\subsection{Dijkstra algorithm}




\subsection{Asymtotic analysis of DFS and BFS}

Memory and time;


what  is branching factor?
if is variable then uses the mean.



how traverse a tree is the question:


traversing is visit each node of graph:



\subsection{question}
whats mean greedy algorithm
what means recursive algorithm



whats is queue?



\section{how represent graphs in python}

\definecolor{xdxdff}{rgb}{0.49019607843137253,0.49019607843137253,1}
\begin{tikzpicture}[line cap=round,line join=round,>=triangle 45,x=1cm,y=1cm]
\clip(4.4460283220320145,3.4839563756562133) rectangle (26.028671106012656,17.35447680358096);
\draw [line width=2pt] (13.9,14.74) circle (1cm);
\draw [line width=2pt] (12,12.16) circle (1cm);
\draw [line width=2pt] (16.08,12.36) circle (1cm);
\draw [line width=2pt] (12.105633182304109,9.257625701340354) circle (1cm);
\draw [line width=2pt] (16.205875233780116,9.200775362612232) circle (1cm);
\draw [line width=2pt] (13.303297233655439,13.937537659053866)-- (12.371390676354103,13.08847669088526);
\draw [line width=2pt] (14.378852130680574,13.862104427085615)-- (15.315088801882949,13.004135745782781);
\draw [line width=2pt] (12.091591730846122,11.164203356683398)-- (12.073162352429183,10.257098384912663);
\draw [line width=2pt] (16.1610200714651,11.363287529916581)-- (16.197271994703605,10.200738354066113);
\draw (13.796623559594307,15.10863897461962) node[anchor=north west] {$\mathbf{A}$};
\draw (11.833281055030623,12.566181055040746) node[anchor=north west] {$\mathbf{B}$};
\draw (16.014211856115878,12.74980301589922) node[anchor=north west] {$\mathbf{C}$};
\draw (11.918029652349919,9.64235444752504) node[anchor=north west] {$\mathbf{D}$};
\draw (16.02833662233576,9.628229681305157) node[anchor=north west] {$\mathbf{E}$};
\draw [line width=2pt] (12.999207410771165,12.199806409785158)-- (15.088738774006357,12.228086460734371);
\draw [line width=2pt] (13.07383916782454,9.507779992933388)-- (15.547220496200314,11.513745900848354);
\begin{scriptsize}
\draw [fill=xdxdff] (13.303297233655439,13.937537659053866) circle (2.5pt);
\draw [fill=xdxdff] (12.371390676354103,13.08847669088526) circle (2.5pt);
\draw [fill=xdxdff] (14.378852130680574,13.862104427085615) circle (2.5pt);
\draw [fill=xdxdff] (15.315088801882949,13.004135745782781) circle (2.5pt);
\draw [fill=xdxdff] (12.091591730846122,11.164203356683398) circle (2.5pt);
\draw [fill=xdxdff] (12.073162352429183,10.257098384912663) circle (2.5pt);
\draw [fill=xdxdff] (16.1610200714651,11.363287529916581) circle (2.5pt);
\draw [fill=xdxdff] (16.197271994703605,10.200738354066113) circle (2.5pt);
\draw [fill=xdxdff] (12.999207410771165,12.199806409785158) circle (2.5pt);
\draw [fill=xdxdff] (15.088738774006357,12.228086460734371) circle (2.5pt);
\draw [fill=xdxdff] (13.07383916782454,9.507779992933388) circle (2.5pt);
\draw [fill=xdxdff] (15.547220496200314,11.513745900848354) circle (2.5pt);
\end{scriptsize}
\end{tikzpicture}

The last figure is a undirected graph;
we have some representations to this graph, by dictionaries or by adjacency matrix.

\begin{lstlisting}
# defining a graph 
graph = {
'A':[B,C]
'B':[A,C]
'C':[D,E]
'E':[C]
}
\end{lstlisting}
Result that we can also represent in a adjacency matrix in the following sense:

\begin{table}[h]
\begin{tabular}{llllll}
 & A & B & C & D & E \\ \cline{2-6} 
\multicolumn{1}{l|}{A} & 0 & 1 & 1 & 0 & 0 \\
\multicolumn{1}{l|}{B} & 1 & 0 & 1 & 1 & 0 \\
\multicolumn{1}{l|}{C} & 1 & 1 & 0 & 1 & 1 \\
\multicolumn{1}{l|}{D} & 0 & 1 & 1 & 0 & 0 \\
\multicolumn{1}{l|}{E} & 0 & 0 & 1 & 0 & 0
\end{tabular}
\end{table}
the table put the format, graph is the importance 


\newpage
in directed graphs (only one direction is allowed) the representation is similar.
The algorithms rely on in data structures,






\subsection{Konigsbiger bridge problems}
resolver by euler an analytical solution an a brute force solution.

$V = (A,B,C,D)$
and $E = \{ (A,B), (A,B) , (A,C), (A,C) , (B,D) , (A,D), (C,D)  \}$
note here that appear two times two different edges.

\begin{lstlisting}
graph = {
'A':[B,B,C,C,D]
'B':[A,A,D]
'C':[A,A,D]
'D':[A,B,C]
}
\end{lstlisting}

in adjacency matrix born the concept of weights,thus:


\begin{table}[h]
\begin{tabular}{lllll}
 & A & B & C & D \\ \cline{2-5} 
\multicolumn{1}{l|}{A} & 0 & 2 & 2 & 1 \\
\multicolumn{1}{l|}{B} & 2 & 0 & 0 & 1 \\
\multicolumn{1}{l|}{C} & 2 & 0 & 0 & 1 \\
\multicolumn{1}{l|}{D} & 1 & 1 & 1 & 0
\end{tabular}
\end{table}






\section{Maze-solving algorithms}
perfect mazes are three from stand point of graph theory.

\subsection{random mouse algorithm}



\subsection{wall follower}


\subsection{pledge algorithm}



\subsection{Hill climbing algorithm}

Stuck in a local maximum.



\subsection{branch and bound algorithm}




\section{A* search algorithm}



\section{tabu search}


\subsection{Grey wolf optimizer}




















\section{Neural networks MIT}



\subsection{Deep Learning}



\subsection{What is backpropagation}


\subsection{softmax}




\subsection{Recurrent neural network}




\subsection{LSTM network}

Is important the idea behind threshold!



\subsubsection{Representation of weights and counter map}

follow the gradient...


\subsubsection{Theorem of gradient}

\subsubsection{Whay is the steepest way??}


\subsection{When gradient descent dont work}

discontinuous function:

gradient require a continuous space:














\end{document}